% ---- pg-425 （书页 413） ----

\noindent 9.9\quad 在法国的一项酒类调查中, 斯托塞尔(Stoetzel, 见文献[14])收集了 $n=1\ 442$ 个人对 $p=9$ 种酒的偏好, 并进行排名. 对排名顺序的 $9\times9$ 样本相关系数矩阵作的因子分析给出如下的载荷估计:

\begin{mvtable}
\small
\begin{tabular}{lccc}
\toprule
变量$(X_i)$ & \multicolumn{3}{c}{因子载荷估计}\\
\cmidrule(lr){2-4}
& $F_1$ & $F_2$ & $F_3$\\
\midrule
(蒸馏)酒 & 0.64 & 0.02 & 0.16\\
Kirsch 酒 & 0.50 & $-0.06$ & $-0.10$\\
Mirabelle 酒 & 0.46 & $-0.24$ & $-0.19$\\
朗姆酒 & 0.17 & 0.74 & $0.97^{*}$\\
白兰地 & $-0.29$ & 0.66 & $-0.39$\\
威士忌 & $-0.29$ & $-0.08$ & 0.09\\
Calvados 酒 & $-0.49$ & 0.20 & $-0.04$\\
科涅克白兰地 & $-0.52$ & $-0.03$ & 0.42\\
Armagnac 酒 & $-0.60$ & $-0.17$ & 0.14\\
\bottomrule
\end{tabular}
\end{mvtable}

\noindent{\footnotesize * 这个数据太大. 斯托塞尔采用了估计因子载荷的近似法, 结果, 它超过了最大值 0.64.}

已知这些结果, 斯托塞尔结论如下: 在法国关于酒的偏好最主要的原则是甜酒和烈酒之间的区别. 第二个激励因素是价格, 记住酒既昂贵又是种摆阔性的消费品, 就可理解这一点. 除了两种最大众化和最便宜的酒(朗姆酒和白兰地), 在产生对酒的喜好的判断上, 这第二个因子起的作用小得多, 第三个因子涉及社会性和地区性的判断差异(见文献[14]).

(a) 就你对所涉及的各种酒的知识来判断, 斯托塞尔的解释合理吗?

(b) 作图表示前两个因子的载荷对. 用图解法表示因子轴的正交旋转. 产生近似旋转载荷. 解释前两个因子的旋转载荷. 你的解释与斯托塞尔从未旋转载荷得到的因子解释一致吗? 请说明.

\noindent 9.10\quad 对鸡的骨骼测量值(见例 9.14)的相关系数矩阵为
\[
\left[\begin{array}{rrrrrr}
1.000& & & & & \\
0.505&1.000& & & & \\
0.569&0.422&1.000& & & \\
0.602&0.467&0.926&1.000& & \\
0.621&0.482&0.877&0.874&1.000& \\
0.603&0.450&0.878&0.894&0.937&1.000
\end{array}\right]
\]
以下因子载荷估计是由极大似然法求出的:

\begin{mvtable}
\small
\begin{tabular}{lcc|cc}
\toprule
& \multicolumn{2}{c}{因子载荷估计} & \multicolumn{2}{c}{因子载荷的最大方差旋转估计}\\
\cmidrule(r){2-3}\cmidrule(l){4-5}
变\quad 量 & $F_2^{*}$ & $F_1$ & $F_2$ & $F_1^{*}$\\
\midrule
1. 颅骨长度 & 0.602 & 0.200 & 0.484 & 0.411\\
2. 颅骨幅宽 & 0.467 & 0.154 & 0.375 & 0.319\\
3. 股骨长 & 0.926 & 0.143 & 0.603 & 0.717\\
4. 胫骨长 & 1.000 & 0.000 & 0.519 & 0.855\\
5. 肱骨长 & 0.874 & 0.476 & 0.861 & 0.499\\
6. 尺骨长 & 0.894 & 0.327 & 0.744 & 0.594\\
\bottomrule
\end{tabular}
\end{mvtable}
